\section{Conclusion}
\label{sec:conclusion}

Constructing representative microdata for sub-national policy
analysis is two estimation problems, not one: filling in what the
base survey does not observe, and choosing --- then weighting --- the
records that will represent the population at every geography a
policymaker asks about. This paper presented both stages as they are
implemented and deployed in \populace{}: a weighted, regime-gated,
chained quantile-regression-forest imputation operator whose design
choices are each attributed by ablation under a population-view
harness, and a target-informed $L_0$ selection-and-calibration step
whose 57{,}240-record support, refit after selection, outperforms
both dense calibration and random subsampling on a 32{,}633-target
administrative surface and ships as the certified United States
release. A companion incidence analysis of the One Big Beautiful
Bill Act consumes that release end to end, making the pipeline's
choices --- and its disclosed defects --- auditable in a live
legislative application at national, state, and district grain.

The near-term agenda follows directly: promote district target
families into the certified gate; map the penalty--concentration
frontier and extend the dense-calibration budget; add
multiple-imputation variance and principled tail extension to the
fill stage; and grow the population-view harness into a standing
release gate, scoring every candidate build against every survey
view and administrative surface before it ships. The estimator is
installable independently of the stack as \populacefit{}, and every
number in both stages reproduces from committed runs in the
accompanying repositories.
